% Appendix, from exam/paper_b_solutions.md: the model answers to Paper B, with marks.

% The answers' headings are labelled ansb:q:N.
\renewcommand{\answerprefix}{ansb}

\chapter{Model Answers to Paper B}
\label{ansb:ch:answers}

Marks are shown per part. Method carries them: a correct diagnosis with a clumsy fix is worth more
than a correct fix with no diagnosis, and several later parts consume earlier answers, so an error
should be followed through rather than penalised twice.

Code is marked on semantics. Missing semicolons, missing includes and misplaced braces cost nothing.
The questions are in Appendix~\ref{exb:ch:paper}.

% ==========================================================================================
\answersection{1}{The vocabulary}{12}

\answerpart{a}{3}
\textbf{To the compiler, none.} \code{#include} behaves identically whatever the extension, and the
extension of a \emph{header} has no bearing on the language it is compiled as. What decides that is
the compiler invocation --- \code{gcc} against \code{g++} --- or the extension of the
\textbf{source} file that pulled the header in. \worth{1}

\textbf{By convention:} \code{.h} is shared ground. A \code{.h} file may be valid C, valid C++, or
both, and the extension alone does not say which. \code{.hpp} states that the header uses C++-only
constructs --- classes, templates, namespaces, references --- and therefore needs a C++ compiler.
\worth{1}

\textbf{The vendor's register header carries \code{.h}.} It is shipped for C users and included from
C++ applications through an \code{extern "C"} block, so it must compile as either. \code{.hpp} would
claim it is C++-only, which is exactly the thing it must not be. \worth{1}

\answerpart{b}{3}
\textbf{An anonymous namespace} restricts everything declared inside it --- variables, functions,
types and compile-time constants --- to a single translation unit. It replaces C's \code{static} on
file-scope functions and variables, and it is used in \code{.cpp} files to hide internal helpers.
\worth{1}

\textbf{The purpose of \code{driver::gpio}} is to avoid symbol collisions and to mirror the
structure of the project, so that \code{driver::gpio::init()} and \code{driver::timer::init()} can
both exist and the reader can see at the call site which module a function belongs to. \worth{1}

\begin{cppcode}
// Pre-C++17.
namespace driver
{
namespace gpio
{
void init() {}
} // namespace gpio
} // namespace driver

// Compact form, C++17 and later.
namespace driver::gpio
{
void init() {}
} // namespace driver::gpio
\end{cppcode}
\noindent\worth{1}

\answerpart{c}{3}
\textbf{Three advantages}, any three of: a reference cannot be null; the call site needs no
\code{&}; reading it needs no \code{*}; the syntax at the point of use is that of an ordinary
variable, so the code reads the same whether the argument is a reference or a value. \worth{2}

\textbf{What a pointer can do that a reference cannot:} be \emph{reseated}. A reference is bound to
one object when it is initialized and can never be made to refer to another. A pointer can be
reassigned, and can be null to mean ``nothing''.

\textbf{Where it matters in the book:} \code{container::Vector<T>} holds its storage as
\code{T* myData}, because \code{resize()} and \code{pushBack()} call \code{std::realloc} and the
block's address changes --- a reference could not follow it. The raw-pointer factory in Chapter~4
is the other case: it must hand back the address of a heap object, and only a pointer can carry
that. \worth{1}

\answerpart{d}{3}
\begin{cppcode}
((reg |= (static_cast<T>(1U) << bits)), ...);
\end{cppcode}
\noindent\worth{2}

{\itshape Award 1 mark for the loop form \code{for (const auto& bit : {bits...})}, which is correct
but is not what the question asked for.\par}

\medskip
{\itshape A note for the marker. The operand of a comma fold must be a cast-expression, so the
assignment needs its own pair of parentheses; written as
\code{(reg |= (static_cast<T>(1U) << bits), ...)} --- one pair short, the form
\secref{c1:sec:templates} warns against --- GCC rejects it with ``binary expression in operand of
fold-expression''. \textbf{Award full marks either way.} The paper marks semantics, and the missing
parentheses are exactly the class of slip the rubric says costs nothing.\par}

\medskip
\textbf{What the compiler generates.} The pack is expanded and the operation performed once per
element, folded together with the comma operator:

\begin{cppcode}
reg |= (static_cast<T>(1U) << 1U);
reg |= (static_cast<T>(1U) << 2U);
reg |= (static_cast<T>(1U) << 3U);
\end{cppcode}

\noindent This is one instantiation of \code{set}, not three --- but the pack's length and element
types are part of the template arguments, so a call with a \emph{different number} of bits
instantiates a different specialization. That is the code-size caveat for parameter packs.
\worth{1}

% ==========================================================================================
\answersection{2}{The same driver, twice}{12}

\answerpart{a}{6}
\begin{cppcode}
/**
 * @file PWM driver implementation.
 */
#pragma once

#include <cstdint>

namespace driver::pwm
{
/**
 * @brief PWM driver.
 */
class Pwm final
{
public:
    explicit Pwm(const std::uint8_t pin, const std::uint16_t period_ms = 20U) noexcept
        : myPin{pin}
        , myPeriod_ms{period_ms}
        , myEnabled{false}
    {}

    ~Pwm() noexcept = default;

    [[nodiscard]] std::uint8_t pin() const noexcept { return myPin; }
    [[nodiscard]] std::uint16_t period_ms() const noexcept { return myPeriod_ms; }
    [[nodiscard]] bool isEnabled() const noexcept { return myEnabled; }

    void setEnabled(const bool enabled) noexcept { myEnabled = enabled; }

    Pwm()                      = delete; // No default constructor.
    Pwm(const Pwm&)            = delete; // No copy constructor.
    Pwm(Pwm&&)                 = delete; // No move constructor.
    Pwm& operator=(const Pwm&) = delete; // No copy assignment.
    Pwm& operator=(Pwm&&)      = delete; // No move assignment.

private:
    /** Pin the PWM output is connected to. */
    const std::uint8_t myPin;

    /** PWM period in milliseconds. */
    const std::uint16_t myPeriod_ms;

    /** Indicate whether the PWM output is enabled. */
    bool myEnabled;
};
} // namespace driver::pwm
\end{cppcode}

{\itshape Marks, 1 each, capped at 6: \code{namespace driver::pwm} and \code{class Pwm final};
\code{explicit} constructor with the default argument \code{= 20U}; \code{const} on \code{myPin}
and \code{myPeriod_ms}, both private with the \code{my} prefix; \code{[[nodiscard]]} on the three
queries; trailing \code{const} on the three queries; \code{noexcept} throughout; the deleted copy
and move operations.\par}

\medskip
{\itshape Do not deduct for \code{#pragma once} or the includes; do deduct if a default argument
appears anywhere other than the declaration.\par}

\answerpart{b}{4}
{\itshape (1 mark each, any four.)\par}
\begin{enumerate}
  \item \textbf{The null \code{self}.} Every C function has to test it, the test costs a branch on
    every call, and it has to be remembered in each new function somebody adds. A member function is
    called \emph{through} an object; \code{backlight.setEnabled(true)} cannot be written without a
    \code{backlight}. The check is not optimized away --- it ceases to have a case to check.
  \item \textbf{A forgotten \code{pwm_init}.} \code{pwm_t backlight;} is a perfectly legal C object
    full of whatever was on the stack, and \code{pwm_set_enabled(&backlight, true)} will happily
    operate on it. A constructor cannot be skipped: the object does not exist until it has run.
  \item \textbf{An error returned as data.} \code{pwm_is_enabled(NULL)} returns \code{false}, which
    is indistinguishable from a genuine ``not enabled''. The caller cannot tell a failure from an
    answer. The C++ query has no failure mode to disguise.
  \item \textbf{Direct writes to the members.} \code{backlight.pin = 12;} compiles in C, and from
    then on the driver silently drives the wrong pin. Making the members private puts them out of
    reach, and \code{const} on \code{myPin} means even the class cannot change it after
    construction.
  \item \textbf{Accidental duplication.} \code{pwm_t copy = backlight;} gives two objects for one
    peripheral. The deleted copy operations turn that into a compile error.
\end{enumerate}

\answerpart{c}{2}
\textbf{Your header is \code{.hpp}.} It contains a class, \code{noexcept}, default arguments and
deleted operations --- none of which a C compiler can read --- so the extension should say so.
\worth{1}

\textbf{The vendor's header stays \code{.h}}, because it must remain compilable as C. Before your C++
driver can call into that library, its declarations need a guarded \code{extern "C"} block:

\begin{ccode}
#ifdef __cplusplus
extern "C" {
#endif
/* declarations */
#ifdef __cplusplus
}
#endif
\end{ccode}

\noindent Without it the C++ compiler mangles the names and the linker cannot match them against the
C library's unmangled symbols. \worth{1}

% ==========================================================================================
\answersection{3}{Enumerations, files, and lifetime}{12}

\answerpart{a}{5}
\begin{cppcode}
namespace driver::gpio
{
/**
 * @brief GPIO direction configuration.
 */
enum class Direction : std::uint8_t
{
    Input,       ///< Input without pull-up.
    InputPullup, ///< Input with pull-up enabled.
    Output,      ///< Output.
    Count,       ///< Number of supported directions.
};

[[nodiscard]] constexpr bool isDirectionValid(const Direction direction) noexcept
{
    return static_cast<std::uint8_t>(Direction::Count) > static_cast<std::uint8_t>(direction);
}
} // namespace driver::gpio
\end{cppcode}

{\itshape (2 marks: 1 for the enumeration with a fixed underlying type, 1 for the check.
\code{Count} is the ``whatever is needed'' the question asked for --- the check has nothing to
compare against without it, and it must be last.)\par}

\medskip
\textbf{Three reasons for an enumeration class}, any three of: \worth{2}
\begin{itemize}
  \item \textbf{The enumerators are scoped.} They are named \code{Direction::Input}, not
    \code{Input}, so they cannot collide with anything and need no \code{GPIO_DIRECTION_} prefix.
    The prefix convention that C forces on you is built into the language here.
  \item \textbf{No implicit conversion to an integer.} A plain \code{enum} decays to \code{int} and
    will silently take part in arithmetic, in comparisons against an unrelated enumeration, and in
    overload resolution. An \code{enum class} requires an explicit \code{static_cast}, which is why
    the validity check above has two of them, and why it is obvious when a conversion is happening.
  \item \textbf{The enumerators do not leak into the enclosing scope}, so the namespace stays clean.
  \item \textbf{The underlying type is always fixed} --- \code{int} unless you name another --- so
    the size is known and the enumeration can be forward-declared. A plain \code{enum} gets that
    only if you write the type out, which C++11 allows and almost nobody does.
\end{itemize}

\textbf{What fixing the underlying type buys.} The size is known, minimal and the same on every
compiler --- one byte instead of whatever \code{int} happens to be --- which matters when the value
is a member of a struct that goes into a frame or is stored in RAM in quantity. Choosing an
\emph{unsigned} type also means no negative value can exist, which reduces validation to the single
upper-bound comparison above. \worth{1}

\answerpart{b}{4}
{\itshape (0.5 marks for each of the seven, counting \code{= default} and \code{= delete}
separately, plus 0.5 for a coherent statement of the principle.)\par}

\begin{booktable}{@{}>{\raggedright\arraybackslash}p{7.8em}>{\raggedright\arraybackslash}p{5.6em}L@{}}
  \thead{Construct} & \thead{Where} & \thead{Why} \\
  \midrule
  \code{explicit} & Header~only & Belongs to the declaration; repeating it on the out-of-class
    definition is an error. \\
  Default argument & Header~only & May be specified once per scope; repeating it in the definition
    is an error. \\
  \code{[[nodiscard]]} & Header~only & It is the declaration callers see. Repeating it is legal but
    redundant; the book's rule is declaration only. \\
  Trailing \code{const} & \textbf{Both} & It is part of the function's type. Omit it in the source
    file and you are defining a \emph{different} function, which does not match anything
    declared. \\
  \code{noexcept} & \textbf{Both} & Every declaration must carry the same exception specification;
    omit it in the source file and the definition is rejected for having a different one. Since
    C++17 it is also part of the function type. \\
  \mbox{\code{= default}},\newline \mbox{\code{= delete}} & Header~only & Written in the class,
    they are complete definitions with no body to place anywhere else; \mbox{\code{= delete}} must
    be on the first declaration. \\
\end{booktable}

\textbf{The principle:} anything that is part of the function's \textbf{type} must appear in both
places, or the definition does not define the function that was declared. Anything that describes
how the function may be \emph{called} --- \code{explicit}, default arguments, \code{[[nodiscard]]}
--- belongs to the declaration alone.

\medskip
{\itshape Also creditable, from the same section: \code{const} on a parameter passed \textbf{by
value} is dropped from the header declaration and kept in the definition, because it constrains
only the function body and is not part of the type.\par}

\answerpart{c}{3}
{\itshape (0.5 marks each for the three pairs, 1.5 for the hazard.)\par}
\begin{enumerate}
  \item \textbf{A local object.} Constructed where control reaches its declaration; destroyed at the
    closing brace of its scope, in reverse order of construction.
  \item \textbf{An object at namespace scope.} Constructed during program start-up,
    \textbf{before \code{main} is entered}; destroyed \textbf{after \code{main} returns}, in reverse
    order of construction.
  \item \textbf{An object created with \code{new}.} Constructed as part of the \code{new}
    expression; destroyed only when \code{delete} is called on the pointer. If \code{delete} never
    happens, the destructor never runs --- that, and not merely the unreleased memory, is what a
    leak costs you.
\end{enumerate}

\textbf{The hazard with namespace-scope objects.} Their constructors run before \code{main}, and
therefore before any clock, power or peripheral initialization the program performs in \code{main}.
A driver object at namespace scope may configure a peripheral that is not yet clocked. On top of
that, the order of construction between objects in \emph{different} translation units is
unspecified, so one such object cannot safely depend on another. And because they are destroyed
after \code{main} returns, anything still running at that point --- a detached thread, for example
--- is touching objects whose lifetime is ending.

% ==========================================================================================
\answersection{4}{What the compiler writes for you}{13}

\answerpart{a}{4}
\textbf{Why the copy constructor works.} A \code{const} non-static data member can be
\textbf{initialized}, once, in the member initializer list, at the moment the object comes into
existence. A copy constructor is building a brand-new object, so \code{myPin{other.myPin}} is a
legal initialization and nothing is being modified. \worth{1.5}

\textbf{Why the copy assignment operator cannot be written.} Assignment operates on an object that
already exists and whose \code{myPin} was fixed when it was constructed. The operator would have to
write \code{myPin = other.myPin;}, which is ill-formed --- a \code{const} object is never
assignable. There is no ordering of statements that gets round this: the member simply cannot be
made to hold a different value for the rest of its life. The simplified operator in
\secref{c2:app:b} compiles precisely because it copies only \code{myState}: the target keeps its own
pin, so the result is not a copy. \worth{1.5}

\textbf{What the compiler does.} It still implicitly \emph{declares} a copy assignment operator, but
because the class has a \code{const} non-static data member, that operator is \textbf{defined as
deleted}. \code{led2 = led1;} is therefore a compile error about a \emph{deleted function} rather
than an assignment to a \code{const}: GCC reports
\code{use of deleted function 'Gpio& Gpio::operator=(const Gpio&)'}, and only its follow-up note
names the \code{const} member, which is why the message can be puzzling if you have not met the
rule. The same applies to the move assignment operator. The class in \secref{c2:app:b} deletes all
four explicitly anyway, which states the intent and produces a clearer diagnostic. \worth{1}

\answerpart{b}{4}
\begin{cppcode}
Buffer(Buffer&& other) noexcept
    : myData{other.myData}
    , mySize{other.mySize}
{
    other.myData = nullptr;
    other.mySize = 0U;
}
\end{cppcode}

{\itshape (2 marks: 1 for taking the members, 1 for clearing \code{other}.)\par}

\medskip
Clearing \code{other.myData} is not tidiness. \code{other} is still a live object and its
destructor will run and \code{delete[]} whatever it holds; leaving the old address there means the
block this object now owns is freed behind its back. \code{delete[] nullptr} is defined and does
nothing, so nulling the pointer makes the moved-from object harmless while keeping it destructible
--- which is the whole contract a moved-from object owes.

\textbf{What move assignment must do that a move constructor need not:} release the resource the
target \textbf{already owns} before taking the new one, and guard against self-assignment.
\worth{1}

\textbf{Why the constructor is exempt.} It is building an object that did not exist a moment ago.
It owns nothing yet, so there is nothing to free; and the object under construction cannot possibly
be the same object as \code{other}, so \code{this != &other} is not merely unnecessary, it can
never be false. \worth{1}

\answerpart{c}{5}
\textbf{\code{= default}} asks the compiler for its own implementation of a special member
function, stated explicitly. The reader can see that the function exists and was intended, and you
can attach \code{noexcept}, \code{virtual} or an access level to it while you are there.

\textbf{\code{= delete}} removes the function entirely. Any attempt to use it is a compile error
that names the deleted function, rather than a program that compiles and does something you did not
want. \worth{2}

\begin{cppcode}
Gpio()                       = delete; // No default constructor.
Gpio(const Gpio&)            = delete; // No copy constructor.
Gpio(Gpio&&)                 = delete; // No move constructor.
Gpio& operator=(const Gpio&) = delete; // No copy assignment.
Gpio& operator=(Gpio&&)      = delete; // No move assignment.
\end{cppcode}

\noindent\textbf{The design reason:} a \code{Gpio} object stands for one physical pin. A copy would
be a second object that believes it owns the same pin and acts on it independently; nothing in the
hardware was duplicated, only the software's idea of it. Deleting the operations makes the mistake a
build failure. The deleted default constructor is there for a related reason --- a GPIO with no pin
number is not a thing. \worth{1.5}

\textbf{What \code{~Gpio() noexcept = default;} nevertheless changes.} Declaring a destructor at
all --- even defaulted, even in the class body --- makes it a \textbf{user-declared} destructor, and
a class with a user-declared destructor gets \textbf{no implicitly declared move constructor and no
implicitly declared move assignment operator}. The class silently loses its move operations and
falls back to copying wherever a move would have been used. The copy operations are still
generated, so nothing fails to compile and nothing warns.

That is precisely why the book writes all of them out: once you touch one special member function,
the compiler stops supplying the others on the terms you assumed, and the only readable answer is to
state every one of them. \worth{1.5}

% ==========================================================================================
\answersection{5}{Designing against an abstraction}{13}

\answerpart{a}{5}
\begin{cppcode}
/**
 * @file Serial driver interface.
 */
#pragma once

#include <cstdint>

namespace driver::serial
{
/**
 * @brief Serial driver interface.
 */
class Interface
{
public:
    virtual ~Interface() noexcept = default;

    virtual void init(std::uint32_t baudrate) noexcept        = 0;
    virtual void write(const char* message) noexcept          = 0;
    [[nodiscard]] virtual bool isInitialized() const noexcept = 0;
};
} // namespace driver::serial
\end{cppcode}

{\itshape (2 marks for the interface; the trailing \code{const} on the query and \code{= 0} on all
three are the parts worth checking.)\par}

\medskip
\textbf{Four conventions}, 0.75 marks each (\secref{c3:app:b} lists five; the fifth is below):
\begin{enumerate}
  \item \textbf{The destructor is \code{virtual}.} Concrete drivers are deleted through an
    \code{Interface*}. Without \code{virtual} that is undefined behavior, and in practice the
    derived destructor never runs, so whatever the driver acquired is never released.
  \item \textbf{The destructor is \code{= default}.} The interface owns nothing and has nothing to
    clean up, so the compiler's implementation is exactly right. Writing it explicitly shows it was
    decided rather than forgotten --- and it is the only way to state \code{virtual} on it.
  \item \textbf{Every method is pure virtual (\code{virtual ... = 0}).} This is what makes the class
    abstract, so nobody can create an \code{Interface}, and what obliges every concrete driver to
    provide the method. The interface defines \emph{what}, never \emph{how}.
  \item \textbf{\code{noexcept} on every method.} It is part of the function type, and an override
    may not weaken it, so declaring it here forces every implementation on every platform to be
    \code{noexcept}. On an embedded target that is exactly the constraint you want to impose from
    the abstraction downwards.
\end{enumerate}

{\itshape A fifth, worth crediting in place of any of the above: \code{[[nodiscard]]} on the query,
so that a caller who asks whether the port is initialized and then ignores the answer is warned,
where the compiler supports it (\secref{c3:app:b} notes that GCC does not warn for a call made
through the interface itself). Note that the attribute is \textbf{not} inherited --- it belongs to
the declaration it is written on --- so every overriding implementation must repeat it, which is
what the book does.\par}

\answerpart{b}{4}
\textbf{\code{private} against \code{protected}.} Neither is accessible from outside the class. The
difference is that a \textbf{derived} class can reach \code{protected} members and cannot reach
\code{private} ones. That is the entire difference, and it is why a base class intended to be
inherited from puts its shared state and helpers under \code{protected}. \worth{1}

\textbf{The three kinds of inheritance:} \worth{1.5}
\begin{itemize}
  \item \textbf{Public.} Public members of the base stay public in the derived class and protected
    stay protected. Used when the derived class \emph{is a} specialization of the base. By far the
    most common.
  \item \textbf{Protected.} Public and protected members of the base become protected in the derived
    class. Like private inheritance, except that classes derived further down still get access.
  \item \textbf{Private.} Everything from the base becomes private in the derived class. Used when
    the base is an implementation detail of the derived class rather than something it presents to
    the world.
\end{itemize}

\textbf{A concrete driver uses public inheritance}, because the point of the exercise is that a
\code{driver::serial::Esp32s3} \emph{is a} \code{driver::serial::Interface} and can be handed to
anything expecting one. With private or protected inheritance, the conversion from derived to base
is not accessible to outside code, so passing the driver as an \code{Interface&} would not compile
--- the abstraction would exist and be unusable. \worth{1.5}

\answerpart{c}{4}
\textbf{It does not compile.} \code{Interface} has pure virtual members and is therefore an
\textbf{abstract class}. A parameter of abstract class type cannot be declared, because passing by
value would require constructing one, and an abstract class cannot be instantiated. The compiler
says so directly: \emph{``cannot declare parameter `serial' to be of abstract type
`driver::serial::Interface'\,''}, followed by a note listing the pure virtual functions responsible.
\worth{1.5}

\textbf{Had the base not been abstract} --- every method with a default implementation --- it would
have compiled and \textbf{sliced}. Only the base subobject of the argument is copied into the
parameter; the derived part is discarded and the parameter's dynamic type is the base.
\code{serial.write(...)} would then call the base version, not the driver's override. Virtual
dispatch works through references and pointers; a by-value copy has already thrown away the thing
that makes it work. The mechanism is \textbf{object slicing}, and it is the more dangerous of the
two outcomes because it compiles. \worth{1.5}

\textbf{The corrected signature:}

\begin{cppcode}
void sendGreeting(driver::serial::Interface& serial) noexcept
\end{cppcode}

\noindent\textbf{Why not \code{const}.} \code{write()} changes the state of the port and is
therefore not a \code{const} method, so a \code{const Interface&} could not call it. A \code{const}
reference to an interface restricts you to its query methods --- here, \code{isInitialized()} alone.
\worth{1}

% ==========================================================================================
\answersection{6}{Factories}{13}

\answerpart{a}{3}
An arrow reads ``depends on'', that is, includes the header of:

\begin{console}
app::logic::Logic          --> driver::factory::Interface, driver::gpio::Interface
driver::factory::Interface --> driver::gpio::Interface     (returned by gpio())
driver::factory::Esp32s3   --> driver::factory::Interface  (implements)
                           --> driver::gpio::Esp32s3       (creates)
driver::factory::Stub      --> driver::factory::Interface  (implements)
                           --> driver::gpio::Stub          (creates)
driver::gpio::Esp32s3      --> driver::gpio::Interface     (implements)
driver::gpio::Stub         --> driver::gpio::Interface     (implements)
\end{console}

{\itshape (1.5 marks for the four layers and the direction of the arrows. The architecture overview
in \secref{c4:app:a} draws the same four layers as a stack, top to bottom; a candidate who
reproduces it has the layers but has drawn who sits above whom, not what depends on what.)\par}

\medskip
\textbf{The dependency that makes it testable} is the one from \code{Logic} to the
\textbf{interfaces} rather than to any concrete type. Because \code{Logic} names only
\code{driver::factory::Interface} and \code{driver::gpio::Interface}, the entire lower half of the
graph can be replaced without \code{Logic} changing, or even recompiling for a different reason.
Logic points only at abstractions, and nothing points from an abstraction down to an
implementation. That is the property, and swapping in stubs is one consequence of it.
\worth{1.5}

\answerpart{b}{4}
\begin{cppcode}
#pragma once

#include <cstdint>
#include <memory>

#include "driver/factory/interface.hpp"
#include "driver/gpio/stub.hpp"

namespace driver::factory
{
class Stub final : public Interface
{
public:
    Stub() noexcept           = default;
    ~Stub() noexcept override = default;

    [[nodiscard]] std::unique_ptr<gpio::Interface> gpio(const std::uint8_t pin) noexcept override
    {
        // Ignore the pin number since it is not used by the stub.
        (void)(pin);
        return std::make_unique<gpio::Stub>();
    }

    Stub(const Stub&)            = delete;
    Stub(Stub&&)                 = delete;
    Stub& operator=(const Stub&) = delete;
    Stub& operator=(Stub&&)      = delete;
};
} // namespace driver::factory
\end{cppcode}

{\itshape (1 mark each, capped at 4: \code{final} and \code{: public Interface}; \code{override}
on both the method and the destructor; the return type matching the interface exactly and
\code{std::make_unique}; \code{(void)(pin)}; \code{[[nodiscard]]} repeated on the override; the
deleted copy and move operations.)\par}

\medskip
{\itshape On \code{[[nodiscard]]}: the attribute belongs to the declaration it is written on and is
\textbf{not} inherited by an override, so the book repeats it on every overriding method. A
candidate who omits it has lost one of the six available marks, not four --- but a candidate who
omits it and explains that it is inherited from the interface has got the rule backwards and should
be told so.\par}

\medskip
Worth stating: \code{std::make_unique<gpio::Stub>()} produces a
\code{std::unique_ptr<gpio::Stub>}, which converts implicitly to
\code{std::unique_ptr<gpio::Interface>} because \code{Stub} derives from \code{Interface}. That
conversion is what makes the \code{return} statement work, and it is safe only because
\code{gpio::Interface} has a virtual destructor.

\answerpart{c}{4}
{\itshape (0.75 marks each for the four, 1 mark for the \code{std::move} answer.)\par}
\begin{enumerate}
  \item \textbf{Ownership is stated in the type.} \code{std::unique_ptr} says ``you own this,
    exclusively''. A raw pointer says nothing at all; the answer lives in a comment, or nowhere.
  \item \textbf{Release is automatic.} The \code{unique_ptr}'s destructor deletes the object when it
    goes out of scope. The raw pointer needs a matching \code{delete} on every path out, and the
    class holding it needs a destructor written for the purpose.
  \item \textbf{It cannot be copied.} \code{std::unique_ptr}'s copy constructor is deleted, so two
    owners of one object cannot be created by accident --- the double delete that a copied raw
    pointer produces at run time becomes a compile error instead.
  \item \textbf{Transfer is explicit.} Handing ownership on requires \code{std::move}, so the moment
    ownership changes hands is visible in the source. A raw pointer changes owner by plain
    assignment, and nothing marks it.
\end{enumerate}

\textbf{Where \code{std::move} becomes necessary.} Not in the Chapter~4 constructor:
\code{factory.}\allowbreak\code{gpio(pin)} returns a temporary, and initializing a member from a temporary already
moves. It is needed as soon as the source is a \textbf{named} \code{unique_ptr} --- a logic class
taking one as a parameter and storing it, \code{myLed{std::move(led)}}, or reassigning
\code{myLed = std::move(other)}.

\textbf{Afterwards the source holds \code{nullptr}.} It is still a valid object; it simply owns
nothing. Dereferencing it is undefined behavior, which is why a moved-from pointer should not be
used again.

\answerpart{d}{2}
\textbf{The cost: dynamic allocation.} The factory creates its objects with \code{new} or
\code{std::make_unique}, which means a heap on a target that may not want one: fragmentation over a
long uptime, an allocation time that is not deterministic, and a failure mode that terminates the
program (Paper~A, Question~1(d)). It also adds a level of indirection and one more abstraction to
read through. \worth{1}

\textbf{The alternative shown in Chapter~4}: construct the concrete drivers directly, as objects
with automatic or static storage, and pass references into the logic class.

\begin{cppcode}
driver::gpio::Esp32s3 led{ledPin};
driver::gpio::Esp32s3 button{buttonPin};
app::logic::Logic logic{led, button};
\end{cppcode}

\noindent with \code{Logic} holding \code{driver::gpio::Interface&} members. No heap is involved at
all, everything is sized at compile time, and the logic is still written entirely against the
interface. \Secref{c4:app:a} says as much: this is ``very efficient when dynamic memory allocation
should be limited''. What it gives up is choosing the driver type at run time, and the single point
of construction that a factory provides once there are many drivers to make. \worth{1}

% ==========================================================================================
\answersection{7}{A container with its size in the type}{12}

\answerpart{a}{6}
\begin{cppcode}
/**
 * @file Fixed-capacity array implementation.
 */
#pragma once

#include <cstddef>

namespace container
{
template<typename T, std::size_t Size>
class Array final
{
    static_assert(Size > 0U, "Array size must be greater than zero!");

public:
    Array() noexcept
        : myData{}
        , mySize{}
    {}

    [[nodiscard]] bool push(const T& element) noexcept
    {
        // Reject the element if the array is already full.
        if (Size <= mySize) { return false; }
        myData[mySize++] = element;
        return true;
    }

    [[nodiscard]] std::size_t size() const noexcept { return mySize; }
    [[nodiscard]] constexpr std::size_t capacity() const noexcept { return Size; }

    T& operator[](const std::size_t index) noexcept { return myData[index]; }
    const T& operator[](const std::size_t index) const noexcept { return myData[index]; }

private:
    /** Element storage. */
    T myData[Size];

    /** Number of elements currently held. */
    std::size_t mySize;
};
} // namespace container
\end{cppcode}

{\itshape Marks, 1 each: \code{template<typename T, std::size_t Size>} with \code{Size} as a
\textbf{non-type} parameter; the \code{static_assert} with a message; storage as a member array
rather than a pointer, with no allocation anywhere; \code{push} returning \code{false} when full and
marked \code{[[nodiscard]]}; both \code{operator[]} overloads, the second \code{const} and
returning \code{const T&}; \code{size()} and \code{capacity()}, \code{const} and
\code{[[nodiscard]]}.\par}

\medskip
{\itshape Creditable extras, not required: \code{constexpr} on \code{capacity()}, since the answer
is a template argument and is known at compile time; the copy and move operations, which for a
container are a design decision rather than an automatic deletion.\par}

\answerpart{b}{4}
\textbf{Why the definitions cannot live in \file{array.cpp}.} The compiler generates code for
\code{Array<std::uint8_t, 8U>} only where it sees a use with those arguments \textbf{and} the
definition of what to generate. \file{array.cpp} is compiled alone, contains no use, and generates
nothing. The translation unit that does use it sees only a declaration and emits calls to functions
nobody defined. The build therefore fails at \textbf{link} time with \code{undefined reference}, not
at compile time --- which is why the error mentions nothing about templates. \worth{2}

\textbf{Three distinct types.} A non-type template argument is part of the type exactly as a type
argument is, so changing \code{8U} to \code{16U} produces a different type just as surely as
changing \code{std::uint8_t} to \code{std::uint16_t} does. The three are unrelated: no conversion
between them, no common base, and a function taking \code{Array<std::uint8_t, 8U>&} will not accept
\code{b}. \worth{1}

\textbf{For the binary:} every member function that is actually used is generated once per
instantiation, so \code{push}, once it is called on all three, exists three times over. Code size
scales with the number of instantiations, and a non-type parameter makes those very easy to multiply
--- one set per buffer size anybody in the project happens to pick. \worth{1}

\answerpart{c}{2}
\textbf{The advantage: no heap.} The storage is a member array whose size is known at compile time,
so the object lives on the stack or in \code{.bss}. There is no allocation to fail, no
fragmentation, no \code{std::realloc}, and no non-deterministic call on a path with a deadline.
\code{push} cannot fail for want of memory --- only for want of capacity, and it says so by
returning \code{false}. \worth{1}

\textbf{What you give up.} The capacity is part of the type, so \code{Array<T, 8U>} and
\code{Array<T, 16U>} are unrelated and no ordinary function can accept both --- anything wanting to
must become a template itself, which multiplies the instantiations again. And the full capacity is
paid for in RAM whether it is used or not. \worth{1}

% ==========================================================================================
\answersection{8}{Sharing state between threads}{13}

\answerpart{a}{3}
\textbf{The three conditions}, all of which must hold: \worth{1.5}
\begin{enumerate}
  \item Two or more threads access the same memory location.
  \item At least one of those accesses is a \textbf{write}.
  \item Nothing synchronizes them.
\end{enumerate}

\textbf{What the standard says:} the behavior of the \textbf{whole program} is undefined.
\worth{0.5}

\textbf{Why that is stronger than ``the value may be wrong.''} Undefined behavior does not promise
you a wrong value. The compiler is entitled to assume no race exists and to optimize on that basis:
it may cache the variable in a register, hoist a load out of a loop, reorder the two writes, or read
the variable twice and get different answers within one expression. The failure therefore need not
look anything like a torn value, and the code that breaks may be nowhere near the racing access. It
also means the program can pass every test at \code{-O0} and fail in the field at \code{-O2} ---
``it worked when we tested it'' is not evidence of anything. \worth{1}

\answerpart{b}{4}
\textbf{Is the data race gone? Yes.} Every shared access is now an atomic operation, so condition
(3) of part (a) no longer holds and the program has no undefined behavior from this cause.
\worth{1}

\textbf{Is the program correct? No.} Atomicity per \emph{variable} is not atomicity per
\emph{transaction}. The two fields have to change together, and nothing makes the pair of writes
indivisible. \worth{1}

\textbf{An interleaving:} \worth{1}

TX writes \code{data} and then \code{newData}, exactly as it did under the mutex:

\begin{console}
TX: data.store(1); newData.store(true)   <- message 1
RX: newData.load() == true               <- RX sees message 1
RX: data.load() == 1                     <- and consumes it
TX: data.store(2); newData.store(true)   <- message 2 is announced
RX: newData.store(false)                 <- RX clears the flag: message 2 is lost
\end{console}

\noindent RX's test of \code{newData} and its clearing of it are two separate atomic operations, and
TX can run between them. Let TX write \code{data} for message 2 between RX's two loads instead, and
announce it after RX has cleared the flag, and RX prints message 2 twice and message 1 never. Each
individual access is atomic; the check-then-act sequence is not.

\textbf{Two things a mutex provides that a pair of atomics does not:} \worth{1}
\begin{enumerate}
  \item \textbf{Mutual exclusion over a region}, not over a single access. Several statements run
    without another thread interleaving between them, which is exactly what a multi-field update
    needs.
  \item \textbf{Memory synchronization between threads.} Everything written inside a critical
    section is guaranteed visible to the next thread that acquires the same mutex, and not only the
    fields you remembered to make atomic.
\end{enumerate}

\answerpart{c}{3}
\textbf{\code{std::launch::async}} --- the task runs immediately, in a new thread.
\textbf{\code{std::launch::deferred}} --- no thread is created; the task runs synchronously on the
\textbf{calling} thread, the first time \code{get()} or \code{wait()} is called on the future.
\worth{1}

\textbf{With no policy passed}, the implementation is free to choose either, and nothing in the
source says which you got. \worth{0.5}

\textbf{The failure.} Code written as ``start the checksum, print the frame, then collect the
result'' silently becomes fully sequential when the implementation picks \code{deferred}: the frame
is printed first, and only then does the checksum begin --- on the same thread, inside
\code{get()}. Nothing crashes and nothing warns; the concurrency you believed you had is simply
absent, and the measurement that was supposed to show an improvement shows none. Worse, if
\code{get()} is never called, a deferred task never runs at all. The fix is to always pass the
policy explicitly. \worth{1.5}

{\itshape Worth knowing, not required: the destructor of a future returned by \code{std::async}
with the async policy blocks until the task completes, so discarding the future also turns the call
back into a synchronous one.\par}

\answerpart{d}{3}
\textbf{The sequence}, with three threads: \worth{1.5}
\begin{enumerate}
  \item The \textbf{low}-priority thread locks a mutex and begins a long operation.
  \item The \textbf{high}-priority thread wakes, tries to lock the same mutex, and blocks.
  \item The \textbf{medium}-priority thread wakes and preempts the low-priority thread. It does not
    use the mutex; it is doing entirely unrelated work. The scheduler simply gives the CPU to the
    higher-priority runnable thread, which is a normal and correct scheduling decision.
  \item The medium-priority thread runs on, so the low-priority thread cannot finish and cannot
    release the lock.
  \item The high-priority thread stays blocked --- not on its own work, and not on anything the
    medium thread holds, but because the medium thread is delaying the release.
\end{enumerate}

The high-priority task ends up waiting on the medium-priority one. The intended scheduling order has
been inverted, and the highest-priority thread --- often the one with the deadline, resetting the
watchdog --- is the one that misses it.

\textbf{Why \code{std::mutex} cannot prevent it.} It has no \textbf{priority inheritance}: no
mechanism to temporarily raise the lock holder's priority to that of the highest-priority thread
waiting on it, so that the holder can finish and release before anything of middling priority cuts
in. The preemption in step 3 is a scheduler decision the mutex plays no part in. \code{std::thread}
compounds this by exposing no notion of priority at all --- setting one means going through
\code{native_handle()} and the platform API --- so the standard library gives you nothing to work
with here. \worth{1}

\textbf{What does fix it:} the platform's priority-inheritance mutex --- an RTOS mutex configured
for priority inheritance, or a POSIX mutex with the \code{PTHREAD_PRIO_INHERIT} protocol. If you
are mixing thread priorities on an RTOS, that is what belongs in the code; \code{std::thread}
itself is not the problem. \worth{0.5}
